\section{Discussion}
\label{sec:discussion}

\paragraph{What ``knows more than it shows'' can mean.} In the information sense the claim is always true. A diffusion model computes its clean-sample estimate from its activations and its other inputs, so these hold at least as much as the estimate shows. The claim worth testing is narrower: that the activations hold usable information about the final sample that a readout of the model's own estimate, equally flexible and equally well tuned, cannot extract. We expected to find such information and found almost none. The layers did look ahead under the standard comparison, and the advantage went away once the estimate was read with the same care, apart from a small gain for floating in the smaller main data sets. What remains is the fall rate of jumps. The layers carry it early in denoising and the estimate shows it only later; on the 50-step model the gain survives every control we applied, and on every model it is small and short-lived.

In practical terms, the floating motions and the balance margins of MDM's samples are as visible in its own estimate after a fifth of the denoising steps as in its layers. Corrections and filters can act on that estimate directly. How fast a jump falls is the one place where the layers have more to offer, and only early in denoising.

\paragraph{Recommendations for probing studies.} Our results suggest five checks before concluding that a generator knows more than it shows:
\begin{enumerate}
    \item Read the model's own clean-sample estimate at the same step as carefully as the activations: the same readout family, tuned the same way. Report what the layer adds to the estimate, for every layer and with an estimator that can fall back to the estimate alone, and run a positive control on the score the conclusion rests on. Near its ceiling, a pooled AUC can miss a real gain.
    \item Score within the groups that share a prompt or a scene, and within prompt and frame where timing is shared, so that recognizing the condition cannot carry the result. Report how the comparison changes with the amount of training data: with a few hundred examples, a high-dimensional probe and a one-number baseline can trade places.
    \item Before concluding that a property is absent from the input, try simple fixed transforms of it as well as linear probes. In a video latent where every linear readout is at chance, the magnitude of local frame-to-frame change exposes plausibility.
    \item Look at what a target's variance is made of. A gravity score can mostly measure whether a body floats or falls, which the estimate shows early, instead of how fast it falls, which it shows only later.
    \item For early pruning, compare with scoring the decoded preview, and report the saving a perfect detector would give.
\end{enumerate}

\paragraph{Limitations.} We study kinematic generators: two transformer diffusion models trained on HumanML3D, MDM at two checkpoints and MotionDiffuse. Models with other training data or designs, such as Kimodo \citep{rempe2026kimodo} or tokenized generators, may behave differently. Physics-based generators are outside the question, since gravity and contact hold in them by construction and our labels cannot fail. What the kinematic planner of a hybrid represents is a separate question (below). Labels are computed from joints with heuristic thresholds (contact heights, the 15\,cm support rule, the 2\,m/s$^2$ floating cut-off). On real motion capture they flag almost no floating and recover the fall rate of real jumps with a spread of about 0.7 to 1.0\,m/s$^2$ (Appendix~\ref{app:realmotion}), but we checked one dataset of scripted actions, and the balance margins are probed as continuous quantities rather than failure labels. Our conclusion that the layers add nothing is relative to the readout families we tried, ridge and tuned MLPs on one frame at a time and a small temporal network. A stronger readout could change the comparison. The data-processing argument does not protect this conclusion: it rules out that the estimate holds information the activations and noisy input lack, but allows the reverse, so it is the layer side that a stronger readout could still raise. The fall-rate result comes from jump-heavy runs, and we do not know which computation in the layers carries it, or whether steering along it would change how fast the generated jumps fall. On MDM our probes see only the conditional pass of classifier-free guidance. The early-warning study is simulated from recorded samples. Each decision uses exactly what an online run would see, but retries draw from a finite pool of 12 recorded samples per prompt. The video test reads only the clean VAE latent of the IntPhys development set. It tests the evidence for ``absent from the input'', not the transformer's own advantage, and Esmati et al.'s latent baseline is reconstructed from their description because their public code does not include it. We do not test whether the model causally uses the directions we decode.

\paragraph{Future work.} The natural next test is an on-trajectory replication in a video diffusion model, with a physics label measurable from the generated video (for example the fall rate of a tracked object, a quantity that \citet{kneifl2026physics} probe with simulator labels) and the model's decoded preview as the baseline. The same comparison applies to the kinematic planners inside physics-based hybrids \citep{tevet2024closd} and to models that drive humanoid robots \citep{wang2026motionbricks}, where knowing whether the planner's own estimate already reveals an imbalance would tell the downstream controller when to intervene.
